\chapter{L'espace de conception et les interfaces}\label{ch:espace}
\iterbanner{V0}

\section{Pourquoi une grille discrète}
La formulation QUBO travaille sur des bits. On fixe donc dès la phase~1 une grille de niveaux pour $(m,p,t)$, et c'est \emph{cette} grille que partagent les trois méthodes : NSGA-II (variables entières), la boucle QUBO (bits) et l'énumération exhaustive (front exact). Si elles ne travaillaient pas sur la même grille, la comparaison n'aurait pas de sens.

\begin{center}
\begin{tabular}{lccc}
\toprule
Variable & niveaux & valeurs & bits \\
\midrule
cambrure $m$ & 8 & $0,\ 1,\ \dots,\ 7$~\% & 3 \\
position $p$ & 8 & $0{,}20,\ 0{,}25,\ \dots,\ 0{,}55$ & 3 \\
épaisseur $t$ & 8 & $6,\ 8,\ \dots,\ 20$~\% & 3 \\
\bottomrule
\end{tabular}
\end{center}
Cela fait $8^3=512$ profils et 9 bits, donc \textbf{9 qubits}.

\paragraph{L'encodage binaire.} Chaque variable prend $2^k$ niveaux d'indice $i=\sum_{j=0}^{k-1}2^j\,b_j$, bit de poids faible en premier. Les 9 bits sont concaténés dans l'ordre $(m,p,t)$ :
\begin{lstlisting}
def encode(self, idx):                 # (im, ip, it) -> 9 bits
    bits = []
    for i, nb in zip(idx, self.bits_per_var):
        bits += [(i >> j) & 1 for j in range(nb)]
    return np.array(bits, dtype=int)
\end{lstlisting}
Un polynôme quadratique en $(m,p,t)$ reste quadratique en bits : chaque variable est affine en bits, et $b^2=b$. C'est ce qui permet de passer du modèle de substitution au QUBO (cours compagnon, chapitre QUBO).

\section{Le simulateur compté : \code{AeroSimulator}}
La métrique de comparaison entre NSGA-II et la boucle QUBO est le \textbf{nombre de simulations distinctes} nécessaires pour retrouver le front exact. Il faut donc le compter au même endroit pour les deux méthodes :
\begin{lstlisting}
class AeroSimulator:
    def __call__(self, idx):
        idx = tuple(int(i) for i in idx)
        if idx not in self.cache:                 # seule une NOUVELLE simulation compte
            m, p, t = self.space.values(idx)
            r = evaluate_at_cl(...) if self.case.mode == "cl" else evaluate(...)
            self.cache[idx] = r
            self.history.append(idx)
        return self.cache[idx]

    @property
    def n_sims(self):
        return len(self.cache)
\end{lstlisting}
NSGA-II réévalue souvent les mêmes individus d'une génération à l'autre ; sans cache, il serait pénalisé à tort. La méthode \code{constraints(r)} renvoie les contraintes au format de pymoo ($g\le0$) : épaisseur minimale, portance minimale en mode incidence, et validité.

\paragraph{Le point de vol : \code{FlightCase}.} Il regroupe $\Re=3\cdot10^6$, le mode (\code{'alpha'} ou \code{'cl'}), $\alpha=2^\circ$ ou $C_L^\star=0{,}4$, $C_{L,min}=0{,}3$ et $t_{min}=8$~\%. Ces valeurs sont à figer en équipe avant toute comparaison.

\section{La base exhaustive}
\code{generer\_base.py} évalue les 512 profils dans chaque mode, en quelques secondes, et écrit un CSV par mode. C'est à la fois la vérité de référence pour le front exact, une table de consultation pour développer les phases suivantes, et un outil de diagnostic.

\begin{figure}[htbp]
\centering
% figure générée par tools/make_figures.py : base_cartes
\begin{tikzpicture}
\begin{groupplot}[group style={group size=2 by 1,horizontal sep=3.2cm},width=0.36\linewidth,height=5.6cm,
  xlabel={cambrure $m$ (\%)},ylabel={épaisseur $t$ (\%)},xtick={0,...,7},ytick={6,8,...,20},
  enlargelimits=false,axis on top,title style={font=\small},tick label style={font=\scriptsize},
  colormap={bl}{color=(white) color=(cblue!45) color=(cblue!85!black)}]
\nextgroupplot[title={$\alpha=2^\circ$ imposé},colorbar,colorbar style={ylabel={$C_D$ (counts)},width=2.5mm,tick label style={font=\scriptsize}}]
\addplot[matrix plot*,mesh/cols=8,point meta=explicit] coordinates {(0,6) [61.11] (1,6) [61.4] (2,6) [58.78] (3,6) [53] (4,6) [52.2] (5,6) [66.3] (6,6) [68.25] (7,6) [71.72] (0,8) [64.25] (1,8) [64.6] (2,8) [63.92] (3,8) [62.85] (4,8) [67.1] (5,8) [68.91] (6,8) [71.27] (7,8) [74.91] (0,10) [66.61] (1,10) [66.91] (2,10) [66.75] (3,10) [66.72] (4,10) [69.89] (5,10) [71.46] (6,10) [74.32] (7,10) [78.15] (0,12) [68.68] (1,12) [69] (2,12) [69.17] (3,12) [69.72] (4,12) [72.31] (5,12) [73.97] (6,12) [77.35] (7,12) [81.31] (0,14) [70.65] (1,14) [71.02] (2,14) [71.46] (3,14) [72.79] (4,14) [74.63] (5,14) [76.47] (6,14) [80.18] (7,14) [85.14] (0,16) [72.56] (1,16) [73.01] (2,16) [73.6] (3,16) [75.08] (4,16) [76.57] (5,16) [79.46] (6,16) [84.27] (7,16) [89.53] (0,18) [74.37] (1,18) [74.8] (2,18) [75.38] (3,18) [77.05] (4,18) [78.86] (5,18) [83.49] (6,18) [88.57] (7,18) [94.14] (0,20) [75.8] (1,20) [76.25] (2,20) [77.9] (3,20) [80.23] (4,20) [82.94] (5,20) [87.74] (6,20) [93.08] (7,20) [98.99]};

\nextgroupplot[title={$C_L=0{,}4$ imposé},colorbar,colorbar style={ylabel={$C_D$ (counts)},width=2.5mm,tick label style={font=\scriptsize}}]
\addplot[matrix plot*,mesh/cols=8,point meta=explicit] coordinates {(0,6) [63.37] (1,6) [61.6] (2,6) [59.09] (3,6) [61.54] (4,6) [60.08] (5,6) [59.19] (6,6) [59.3] (7,6) [62.39] (0,8) [66.09] (1,8) [65.51] (2,8) [62.85] (3,8) [63.84] (4,8) [62.74] (5,8) [61.62] (6,8) [61.86] (7,8) [63.42] (0,10) [69.3] (1,10) [67.78] (2,10) [66.19] (3,10) [66.04] (4,10) [65.25] (5,10) [63.92] (6,10) [64.5] (7,10) [65.52] (0,12) [71.67] (1,12) [69.66] (2,12) [68.36] (3,12) [68.18] (4,12) [67.63] (5,12) [66.67] (6,12) [67.05] (7,12) [67.15] (0,14) [73.57] (1,14) [71.5] (2,14) [70.39] (3,14) [70.26] (4,14) [69.89] (5,14) [69.15] (6,14) [69.49] (7,14) [69.98] (0,16) [75.24] (1,16) [73.3] (2,16) [72.37] (3,16) [72.29] (4,16) [72.05] (5,16) [71.46] (6,16) [71.83] (7,16) [72.86] (0,18) [76.67] (1,18) [74.96] (2,18) [74.21] (3,18) [74.38] (4,18) [74.03] (5,18) [73.59] (6,18) [73.79] (7,18) [75.68] (0,20) [77.85] (1,20) [76.33] (2,20) [75.64] (3,20) [76.88] (4,20) [75.71] (5,20) [75.21] (6,20) [76.79] (7,20) [78.65]};

\end{groupplot}
\end{tikzpicture}

\caption{$C_D$ (counts) sur la tranche $p=0{,}40$ de la grille, dans les deux modes ($\Re=3\cdot10^6$). À incidence imposée, la traînée croît à la fois avec l'épaisseur et la cambrure. À portance imposée, la cambrure ne coûte presque plus rien : elle remplace l'incidence. Les croix marquent les profils invalides.}\label{fig:base}
\end{figure}

\begin{center}
\begin{tabularx}{\linewidth}{Xcc}
\toprule
 & $\alpha=2^\circ$ imposé & $C_L=0{,}4$ imposé \\
\midrule
profils valides & 512 / 512 & 507 / 512 \\
profils faisables (valides, $t\ge8$~\%, et $C_L\ge0{,}3$ en mode $\alpha$) & 392 & 446 \\
plage de $C_D$ (valides) & 51 à 124 counts & 51 à 88 counts \\
corrélation de $C_D$ avec $t$ / avec $m$ & 0,60 / 0,62 & 0,74 / 0,04 \\
meilleur profil faisable $(m,p,t)$ & (4~\%, 0,35, 8~\%) & (4~\%, 0,55, 8~\%) \\
\quad son $C_D$ & 55 counts & 54 counts \\
incidence & $2^\circ$ pour tous & $-5{,}5^\circ$ à $+3{,}5^\circ$ \\
drapeau \code{te\_separation} & 145 profils & 64 profils \\
\bottomrule
\end{tabularx}
\end{center}

\paragraph{Lecture.} Le choix du mode change la forme du problème :
\begin{itemize}
\item à incidence imposée, on compare des profils qui ne portent pas le même poids. Un profil très cambré « gagne » de la portance qu'il paie en traînée, d'où la corrélation de 0,62 avec $m$. Il faut une contrainte $C_L\ge C_{L,min}$ ;
\item à portance imposée, on compare à mission égale, comme un avionneur. La corrélation avec $m$ disparaît, l'épaisseur domine, et la contrainte de portance disparaît du QUBO. En contrepartie, l'incidence devient une \emph{sortie} de l'aéro et une \emph{entrée} de la SER : l'attitude vue par le radar dépend du profil. C'est un vrai couplage multiphysique.
\end{itemize}
La recommandation de la note d'interfaces est le mode portance imposée, avec le mode incidence comme repli.

\section{Vers la phase 2 : la géométrie pour la SER}
La SER a besoin d'une surface. \code{geometry\_export.extrude} transforme le profil en un tronçon d'aile extrudé selon l'envergure, facetté en triangles, avec les \emph{mêmes points} que le calcul de panneaux (figure~\ref{fig:extrusion}). Chaque panneau 2D devient un rectangle de deux triangles, et les deux flancs sont triangulés en bandes entre points d'extrados et d'intrados de même indice.

\begin{figure}[htbp]
\centering
\begin{tikzpicture}[x={(1cm,0cm)},y={(0cm,1cm)},z={(0.55cm,0.35cm)},scale=1.25]
\def\prof{(4,0,0) .. controls (3,0.25,0) and (1.5,0.45,0) .. (0.6,0.4,0) .. controls (0.1,0.32,0) and (0,0.1,0) .. (0,0,0) .. controls (0,-0.08,0) and (0.3,-0.2,0) .. (1,-0.22,0) .. controls (2,-0.2,0) and (3,-0.08,0) .. (4,0,0)}
\fill[cblue!8] (4,0,0) -- (4,0,4) -- (0,0,4) -- (0,0,0) -- cycle;
\foreach \z in {0,4}{
  \draw[cink,thick] plot[smooth,tension=0.6] coordinates {(4,0,\z) (3,0.2,\z) (2,0.34,\z) (1,0.38,\z) (0.4,0.3,\z) (0.08,0.15,\z) (0,0,\z) (0.1,-0.1,\z) (0.5,-0.18,\z) (1.5,-0.2,\z) (2.5,-0.14,\z) (3.3,-0.06,\z) (4,0,\z)};}
\foreach \x/\y in {4/0,3/0.2,2/0.34,1/0.38,0.4/0.3,0.08/0.15,0/0}{\draw[cmuted,thin] (\x,\y,0) -- (\x,\y,4);}
\foreach \x/\y/\xx/\yy in {3/0.2/4/0,2/0.34/3/0.2,1/0.38/2/0.34}{\draw[cmuted,thin] (\x,\y,0) -- (\xx,\yy,4);}
\draw[arr,cmuted] (0,0,0) -- (1.2,0,0) node[below,font=\small] {$x$ (corde)};
\draw[arr,cmuted] (4,0,0) -- (4,0,5) node[right,font=\small] {$z$ (envergure)};
\node[font=\small,align=center] at (2.6,-0.9,0) {chaque panneau 2D $\to$ un rectangle\\ de deux triangles};
\draw[corange,very thick,-{Stealth[length=2.5mm]}] (-2.2,-0.4,2) -- (-0.2,-0.02,2) node[pos=0,left,font=\small,corange,align=right] {radar de face,\\ tourné de $\alpha$};
\end{tikzpicture}
\caption{L'extrusion 2,5D exportée en STL. L'incidence n'est pas appliquée à la géométrie : elle est écrite dans les métadonnées JSON, et c'est la direction radar qui doit être tournée.}\label{fig:extrusion}
\end{figure}

Conventions du fichier : repère corps ($x$ le long de la corde vers l'aval, $y$ vers l'extrados, $z$ selon l'envergure, origine au bord d'attaque), normales sortantes, triangles dégénérés retirés au bord d'attaque et au bord de fuite fermé. Deux contrôles garantissent une surface fermée et bien orientée : $\sum_k\vect n_kA_k=0$, et le volume calculé par le théorème de la divergence égal à l'aire du profil fois l'envergure. Pour le NACA 2412 avec 200 panneaux, le STL compte 796 triangles.

\begin{piege}
Le facettage est un \emph{paramètre numérique} pour la SER, pas une variable de conception. À haute fréquence, chaque facette plane produit des reflets spéculaires qui n'existent pas sur la forme lisse. Il faut figer \code{n\_half} pour tout le hackathon et vérifier la convergence de la SER avec lui, comme on l'a fait pour $C_D$.
\end{piege}

\section{Vers la phase 4 : la sortie est-elle « QUBO-compatible » ?}
La boucle quantique remplace le simulateur par un modèle quadratique ajusté par moindres carrés, en variables normalisées $z_i=2x_i/7-1\in[-1,1]$ :
\[
\widehat{C_D}(z)=c_0+\sum_ic_iz_i+\sum_{i\le l}c_{il}\,z_iz_l\qquad(10\ \text{coefficients}).
\]
\code{analyse\_quadratique.py} mesure, sur la base exhaustive, la qualité de ce modèle.

\begin{figure}[htbp]
\centering
% figure générée par tools/make_figures.py : quadratique
\begin{tikzpicture}
\begin{axis}[width=0.62\linewidth,height=6cm,xlabel={$C_D$ du modèle aéro (counts)},ylabel={$C_D$ quadratique (counts)},
  legend style={draw=none,font=\scriptsize,at={(0.02,0.98)},anchor=north west},legend cell align=left]
\addplot[only marks,mark=*,mark size=1.1pt,cblue,fill opacity=0.6,draw opacity=0.6] coordinates {(61.11,58.04) (64.25,60.84) (66.61,63.59) (68.68,66.3) (70.65,68.97) (72.56,71.59) (74.37,74.17) (75.8,76.71) (61.11,57.16) (64.25,59.78) (66.61,62.36) (68.68,64.89) (70.65,67.39) (72.56,69.84) (74.37,72.25) (75.8,74.62) (61.11,57.22) (64.25,59.66) (66.61,62.07) (68.68,64.43) (70.65,66.75) (72.56,69.03) (74.37,71.27) (75.8,73.46) (61.11,58.22) (64.25,60.49) (66.61,62.73) (68.68,64.92) (70.65,67.06) (72.56,69.17) (74.37,71.23) (75.8,73.25) (61.11,60.17) (64.25,62.27) (66.61,64.33) (68.68,66.35) (70.65,68.32) (72.56,70.25) (74.37,72.14) (75.8,73.99) (61.11,63.06) (64.25,64.99) (66.61,66.88) (68.68,68.72) (70.65,70.52) (72.56,72.28) (74.37,74) (75.8,75.67) (61.11,66.9) (64.25,68.65) (66.61,70.37) (68.68,72.04) (70.65,73.67) (72.56,75.25) (74.37,76.8) (75.8,78.3) (61.11,71.68) (64.25,73.26) (66.61,74.8) (68.68,76.3) (70.65,77.76) (72.56,79.17) (74.37,80.54) (75.8,81.87) (56.06,59.39) (60.79,62.43) (64.62,65.42) (67.78,68.38) (70.46,71.29) (72.85,74.16) (74.92,76.99) (76.41,79.78) (56.99,57.97) (60.92,60.83) (64.15,63.66) (67.08,66.44) (69.84,69.18) (72.36,71.88) (74.51,74.53) (76.07,77.15) (58.51,57.5) (61.96,60.19) (64.75,62.84) (67.25,65.45) (69.62,68.02) (71.91,70.54) (74.04,73.02) (75.63,75.46) (60,57.97) (63.28,60.49) (65.78,62.97) (68.04,65.4) (70.2,67.8) (72.29,70.15) (74.2,72.45) (75.65,74.72) (61.4,59.38) (64.6,61.73) (66.91,64.04) (69,66.3) (71.02,68.52) (73.01,70.7) (74.8,72.83) (76.25,74.93) (62.8,61.75) (65.88,63.92) (68.01,66.05) (69.96,68.14) (71.88,70.19) (73.78,72.19) (75.48,74.16) (76.93,76.08) (64.25,65.05) (67.04,67.05) (68.94,69.01) (70.7,70.93) (72.46,72.8) (74.23,74.63) (75.84,76.42) (77.34,78.17) (65.6,69.31) (67.83,71.13) (69.48,72.92) (71.14,74.66) (72.83,76.36) (74.55,78.02) (76.1,79.64) (77.67,81.21) (59.66,64.9) (64.32,68.15) (67.98,71.35) (71.07,74.51) (73.73,77.62) (75.83,80.7) (78.4,83.73) (52.62,59.67) (58.87,62.78) (63.49,65.85) (67.13,68.88) (70.28,71.86) (73,74.81) (75.19,77.71) (77.82,80.57) (53.06,58.66) (59.22,61.6) (63.35,64.5) (66.64,67.36) (69.61,70.17) (55.22,58.61) (61.09,61.37) (64.62,64.1) (67.47,66.78) (70.03,69.42) (72.38,72.01) (74.37,74.57) (58.78,59.49) (63.92,62.08) (66.75,64.63) (69.17,67.14) (71.46,69.61) (73.6,72.03) (75.38,74.41) (77.9,76.75) (63.61,61.32) (67.03,63.74) (69.14,66.12) (71.26,68.45) (73.19,70.75) (75.03,73) (76.49,75.2) (79.06,77.37) (67.91,64.1) (69.63,66.34) (70.99,68.55) (72.52,70.71) (74.19,72.83) (75.85,74.91) (77.1,76.94) (79.9,78.93) (70.18,67.82) (71.14,69.89) (72.19,71.92) (73.54,73.91) (75.08,75.86) (76.62,77.76) (77.93,79.62) (80.71,81.44) (77.35,85.29) (80.94,88.57) (51.16,60.13) (58.99,63.15) (64.2,66.11) (67.67,69.04) (53,60.49) (62.85,63.33) (66.72,66.12) (69.72,68.88) (72.79,71.59) (75.08,74.26) (77.05,76.88) (80.23,79.47) (56.46,61.79) (67.66,64.45) (70.35,67.07) (72.34,69.66) (74.32,72.19) (76.15,74.69) (78.19,77.14) (81.28,79.55) (71.04,64.03) (71.33,66.52) (72.45,68.97) (73.99,71.38) (75.71,73.74) (77.31,76.07) (79.56,78.35) (82.58,80.58) (74.39,67.22) (74.06,69.54) (74.6,71.82) (75.79,74.05) (77.27,76.24) (78.63,78.39) (81.12,80.5) (84.1,82.56) (51.78,62.55) (55.1,65.81) (52.2,62.37) (67.1,65.46) (69.89,68.5) (72.31,71.5) (74.63,74.46) (76.57,77.37) (54.35,63.14) (68.71,66.05) (70.92,68.92) (73.13,71.75) (75.3,74.53) (77.12,77.27) (80.28,79.97) (83.6,82.63) (71.42,64.85) (71.78,67.59) (73.25,70.29) (75.08,72.94) (77,75.55) (79,78.12) (82.11,80.64) (85.38,83.13) (78.01,67.51) (76.09,70.08) (76.51,72.6) (77.81,75.08) (79.33,77.51) (81.53,79.91) (84.51,82.26) (87.69,84.57) (66.3,65.15) (68.91,68.48) (71.46,71.77) (73.97,75.01) (66.37,65.39) (68.9,68.54) (71.45,71.66) (73.97,74.73) (76.31,77.76) (79.1,80.74) (82.53,83.69) (86.13,86.59) (69.82,66.57) (71.48,69.55) (73.62,72.49) (75.91,75.39) (78.02,78.25) (81.24,81.06) (84.67,83.83) (88.24,86.56) (78.86,68.69) (76.86,71.5) (77.8,74.27) (79.47,77) (81.59,79.68) (84.66,82.32) (87.97,84.92) (91.49,87.47) (66.81,68.52) (69.6,71.92) (72.4,75.28) (75.14,78.6) (77.8,81.87) (81.41,85.11) (68.46,69.17) (71.11,72.4) (73.91,75.59) (76.66,78.73) (79.79,81.83) (83.48,84.89) (87.28,87.91) (91.24,90.88) (76.08,70.77) (76.52,73.82) (78.49,76.83) (80.81,79.8) (84.13,82.73) (87.69,85.62) (91.43,88.46) (95.35,91.26) (68.18,72.55) (71.1,76.19) (74.07,79.8) (67.88,72.66) (71.13,76.14) (74.42,79.57) (77.71,82.96) (81.7,86.31) (85.79,89.61) (90.02,92.87) (94.37,96.09) (73.02,73.73) (75.61,77.03) (78.71,80.29) (82.39,83.5) (86.38,86.68) (90.52,89.81) (94.78,92.9) (99.18,95.94)};\addlegendentry{transition extrados : Michel}
\addplot[only marks,mark=*,mark size=1.1pt,corange,fill opacity=0.6,draw opacity=0.6] coordinates {(53.81,61.62) (72.31,72.94) (75.28,75.66) (78.85,78.35) (76.45,77.08) (55.45,64.74) (60.94,68.27) (67.16,71.76) (72.15,75.21) (76.45,78.61) (80.16,81.97) (52.97,62.26) (58.8,65.62) (65.12,68.93) (69.85,72.21) (73.88,75.44) (77.41,78.63) (80.88,81.77) (85.57,84.87) (50.53,60.72) (57.23,63.91) (63.54,67.05) (67.98,70.15) (71.76,73.21) (75.13,76.22) (78.28,79.2) (82.79,82.13) (69.91,71.93) (73.1,74.77) (75.94,77.57) (80.17,80.32) (76.74,68.75) (80.15,72.53) (83.58,76.26) (74.61,79.95) (79.39,83.6) (83.22,87.21) (88.67,90.78) (94.19,94.3) (73.14,65.74) (76.13,69.34) (66.65,72.9) (72.39,76.42) (76.86,79.9) (80.4,83.33) (85.36,86.73) (90.5,90.07) (53.59,63.67) (57,67.1) (65.4,70.49) (70.68,73.84) (74.75,77.14) (78.05,80.4) (82.45,83.62) (87.22,86.79) (65,69.02) (69.47,72.2) (73.16,75.33) (76.24,78.41) (80.15,81.46) (84.57,84.46) (78.86,80.24) (82.94,83.07) (79.65,73.65) (83.29,77.68) (87.03,81.66) (90.81,85.59) (94.53,89.49) (87.78,93.34) (93.76,97.15) (99.85,100.9) (76.11,70.11) (79.47,73.96) (82.89,77.77) (86.34,81.53) (89.51,85.25) (84.71,88.93) (90.18,92.57) (95.8,96.16) (72.62,67.51) (75.7,71.19) (78.8,74.82) (81.92,78.41) (84.8,81.96) (82.24,85.47) (87.21,88.93) (92.3,92.35) (68.97,65.86) (71.86,69.36) (74.78,72.82) (77.71,76.24) (80.53,79.61) (83.52,82.95) (87.56,86.24) (91.71,89.49) (76.47,78.21) (79.46,81.37) (83.49,84.49) (87.74,87.57) (82.81,79.45) (86.66,83.71) (90.69,87.94) (94.82,92.12) (98.79,96.26) (104.5,100.4) (110.2,104.4) (106,108.4) (79.07,75.37) (82.68,79.47) (86.37,83.52) (90.12,87.53) (93.6,91.49) (98.69,95.42) (103.7,99.3) (101.7,103.1) (75.68,72.24) (78.99,76.16) (82.37,80.04) (85.77,83.88) (88.9,87.67) (93.47,91.42) (98.16,95.13) (98.07,98.8) (72.2,70.06) (75.28,73.8) (78.41,77.51) (81.56,81.17) (84.45,84.79) (88.72,88.37) (93.16,91.91) (97.78,95.4) (68.25,68.82) (71.27,72.39) (74.32,75.92) (77.35,79.41) (80.18,82.86) (84.27,86.27) (88.57,89.63) (93.08,92.95) (84.16,88.3) (88.8,91.44) (86.36,86.13) (90.41,90.64) (94.65,95.11) (98.91,99.54) (104,103.9) (110.2,108.3) (116.8,112.6) (123.7,116.8) (82.6,81.52) (86.43,85.86) (90.39,90.16) (94.34,94.41) (98.99,98.63) (104.5,102.8) (110.2,106.9) (116.1,111) (78.97,77.86) (82.54,82.03) (86.16,86.15) (89.77,90.23) (94.02,94.27) (99.06,98.27) (104.3,102.2) (109.7,106.1) (75.55,75.14) (78.85,79.14) (82.24,83.09) (85.56,87) (89.52,90.86) (94.17,94.69) (99.01,98.47) (104.1,102.2) (71.72,73.37) (74.91,77.19) (78.15,80.97) (81.31,84.71) (85.14,88.4) (89.53,92.05) (94.14,95.66) (98.99,99.22) (76.61,83.36) (80.49,86.88) (84.82,90.36) (89.35,93.79) (94.21,97.19)};\addlegendentry{transition extrados : décollement laminaire}
\addplot[cmuted,dashed] coordinates {(50.5,50.5) (123.7,123.7)};
\node[font=\small,anchor=south east] at (rel axis cs:0.98,0.03) {$R^2=0.908$};
\end{axis}
\end{tikzpicture}

\caption{Modèle quadratique à 10 coefficients ajusté sur les 512 profils (mode $\alpha=2^\circ$), contre la vraie valeur. Les couleurs distinguent le mode de transition sur l'extrados. Le $R^2$ de 0,91 est limité par la bascule entre transition de Michel et transition par décollement laminaire.}\label{fig:quad}
\end{figure}

\begin{center}
\begin{tabular}{lcc}
\toprule
Sortie & quadratique en $(m,p,t)$ (10 coef.) & QUBO libre sur 9 bits (46 coef.) \\
\midrule
$C_L$ & $R^2=0{,}9998$ & $0{,}9998$ \\
$C_D$ & $R^2=0{,}908$ & $0{,}920$ \\
\bottomrule
\end{tabular}
\end{center}
\begin{itemize}
\item $C_L$ est quasi affine en $(m,\alpha)$ par construction (panneaux linéaires en $\alpha$, théorie des profils minces linéaire en $m$) : un quadratique le représente parfaitement.
\item Les 9~\% de variance de $C_D$ non expliqués sont structurels : c'est la bascule de mode de transition du chapitre~\ref{ch:transition}. Interpoler $x_{tr}$ entre stations n'a fait passer le $R^2$ que de 0,899 à 0,908.
\item Avec 20 profils tirés au hasard, le quadratique classe bien les profils (corrélation de rang de Spearman d'environ 0,89), mais son minimum n'est dans le vrai top 5~\% qu'une fois sur trois environ.
\end{itemize}

\begin{projet}
Conséquence pour la phase~4 : un QUBO global ajusté une seule fois ne suffit pas. Il faut la \textbf{boucle itérative} de la note d'intention, qui réajuste le modèle autour des meilleurs candidats. C'est ce que montre le cours compagnon : la boucle retrouve le vrai minimum de $C_D$ en 25 simulations (médiane sur 50 essais), contre environ 256 en tirant les profils au hasard.
\end{projet}

\exercices
\begin{exercice}
(a) Donner la chaîne des 9 bits du profil d'indices $(i_m,i_p,i_t)=(4,3,1)$, et les valeurs physiques correspondantes. (b) Combien de chaînes de bits distinctes représentent un profil symétrique ($m=0$) ?
\end{exercice}
\begin{exercice}
Un collègue propose de comparer NSGA-II et la boucle QUBO en temps de calcul plutôt qu'en nombre de simulations distinctes. Donner deux raisons pour lesquelles ce serait une mauvaise métrique dans ce hackathon.
\end{exercice}
